\chapter{Especificaciones Relativas a las Operaciones UAS (OpSpecs)}
\label{anexo:opspecs}
\label{inicio_anexo}

Este anexo contiene la \textbf{propuesta} de Especificaciones Relativas a las Operaciones UAS de Skyfilms S.A.S., elaborada con la estructura del formato MAUT-5.0-12-096 de la Aerocivil. Las OpSpecs vigentes serán las emitidas por la Autoridad junto con el Certificado de Operador UAS (CDO-U); los campos de numeración, fecha y firma los diligencia la Aerocivil. Su contenido es coherente con el alcance operacional del Capítulo de Operaciones.

\section{Información del solicitante}
\begin{table}[H]
    \centering
    \renewcommand{\arraystretch}{1.2}
    \begin{tabular}{|p{5.5cm}|p{9cm}|}
        \hline
        \textbf{CDO-U No.} & Por asignar por la Aerocivil \\ \hline
        \textbf{Revisión No. / Fecha} & 1 / Por asignar por la Aerocivil \\ \hline
        \textbf{Razón social} & Skyfilms S.A.S. \\ \hline
        \textbf{NIT} & 902.039.520-1 \\ \hline
        \textbf{Teléfono} & 301 264 6892 \\ \hline
        \textbf{Correo electrónico} & skyfilmsfpv@gmail.com \\ \hline
        \textbf{Contacto Aerocivil (GDMUA)} & 601 770 4386 Ext. 4386 / gdmua@aerocivil.gov.co \\ \hline
    \end{tabular}
\end{table}

\section{Información de las operaciones}
\begin{table}[H]
    \centering
    \renewcommand{\arraystretch}{1.2}
    \begin{tabular}{|p{5.5cm}|p{9cm}|}
        \hline
        \textbf{Marca, modelo y RUAS} & DJI Mavic 4 Pro Creator Combo, RUAS por asignar \\
        & BetaFPV Pavo20 Pro II, RUAS por asignar \\ \hline
        \textbf{Tipo de operación} & [X] Simple captura de imágenes o datos. Las demás opciones del formato no aplican. \\ \hline
        \textbf{Tipo de contacto visual} & [X] VLOS \quad [\ ] EVLOS \quad [\ ] BVLOS \\ \hline
        \textbf{Limitaciones especiales} &
        \begin{minipage}[t]{8.7cm}
        \begin{itemize}\setlength\itemsep{0pt}
            \item Distancia horizontal máxima de 750 m del piloto remoto, en VLOS; Pavo20 Pro II, máximo 250 m.
            \item Altura máxima de 400 ft AGL; en zona urbana, 30 m sobre el obstáculo más alto en un radio de 150 m, salvo autorización especial.
            \item Vuelo nocturno solo con las luces de fábrica de la aeronave encendidas de forma permanente y con observador.
            \item Pavo20 Pro II (FPV): siempre con observador; vuelo nocturno sujeto a la verificación de sus luces según el MO.
            \item Tiempo de vuelo planificado máximo del 80\,\% de la autonomía.
        \end{itemize}
        \end{minipage} \\ \hline
    \end{tabular}
\end{table}

\section{Vuelos especiales y aprobaciones específicas}
\begin{table}[H]
    \centering
    \renewcommand{\arraystretch}{1.2}
    \begin{tabular}{|p{5cm}|c|c|p{5.2cm}|}
        \hline
        \rowcolor[gray]{0.9} \textbf{Vuelo especial} & \textbf{SÍ} & \textbf{NO} & \textbf{Descripción} \\ \hline
        Vuelo nocturno & X & & Captura de imágenes y datos de noche, según el MO y el RAC 100. \\ \hline
        Vuelo en zona urbana & X & & Captura de imágenes y datos en áreas urbanas, según el MO y el RAC 100. \\ \hline
        Vuelo autónomo & & X & \\ \hline
        Demostraciones comerciales o de capacidad tecnológica & & X & \\ \hline
        Competencias y actividades deportivas o recreativas & & X & \\ \hline
        Vuelos de UA cautiva & & X & \\ \hline
        Otro: espacios cerrados o confinados & & X & No se solicita; no hace parte del alcance del MO. \\ \hline
        Mercancías peligrosas & & X & \\ \hline
        Servicios aéreos comerciales & X & & Producción audiovisual con UAS para producciones, eventos y medios de comunicación. \\ \hline
    \end{tabular}
\end{table}

\section{Gestión del mantenimiento de los UAS}
Skyfilms gestionará el mantenimiento de sus UAS mediante inspecciones, limpieza, ajustes, reemplazo de componentes y reparaciones, de acuerdo con los procedimientos del fabricante y su MCM. Los trabajos que requieran capacidades o autorizaciones específicas serán realizados por personal o talleres técnicamente competentes.

\section{Control de revisiones}
\begin{table}[H]
    \centering
    \renewcommand{\arraystretch}{1.2}
    \begin{tabular}{|c|c|p{5cm}|p{4cm}|}
        \hline
        \rowcolor[gray]{0.9} \textbf{Revisión} & \textbf{Fecha (dd/mm/aaaa)} & \textbf{Razón de la revisión} & \textbf{Revisado por} \\ \hline
        1 & Por asignar & Emisión inicial & Aerocivil \\ \hline
    \end{tabular}
\end{table}
